\documentclass[12pt]{article}
\usepackage[utf8]{inputenc}
\usepackage[T1]{fontenc}
\usepackage{hyperref}
\title{Sani Language Resources for K-12 Curriculum Development}
\author{Compiled by Claude Sonnet 5.5 (claude-sonnet-5-5) for Andr\'as Kornai}
\date{}
\begin{document}
\maketitle

\section{Introduction}
Sani (ISO 639-3: ysn) is a Loloish (Lolo-Burmese, Sino-Tibetan) language of Yunnan, China. It is one of the languages grouped under Southeastern Yi, one of six Yi languages recognised by the Chinese government. It is spoken mainly in Shilin Yi Autonomous County, with speakers also in Luliang, Luxi, Shizong, Yiliang, Malong, Luquan and Mile counties. Wikipedia gives about 100,000 native speakers (2007), and Asia Harvest gave 103,000 in 1999. Sani has its own Yi-derived script, and Omniglot notes it can also be written in the Latin alphabet.

\section{Primary Educational Resources}
No primers, textbooks or school programmes for Sani were found. The language is not known to be taught in schools.

\section{Linguistic and Reference Materials}
\begin{itemize}
\item Wikipedia, ``Sani language''. \url{https://en.wikipedia.org/wiki/Sani_language}
\item Omniglot, ``Sani (Nip-dop)'': writing system and overview. \url{https://www.omniglot.com/writing/sani.htm}
\item BabelStone, ``Sani Script'': notes on the script and its Unicode status. The search summary says the script is related to but not unifiable with the Nuosu Yi script, and that no proposal to encode it has been made. \url{https://www.babelstone.co.uk/Yi/Sani.html}
\item Kobe City University of Foreign Studies annual report paper (nenpo61), which includes a discussion of Sani Yi. \url{https://kobe-cufs.repo.nii.ac.jp/record/2583/files/nenpo61-04.pdf}
\item Kyoto University repository paper in the Sino-Tibetan languages series. Its topic could not be confirmed from the search result. \url{https://repository.kulib.kyoto-u.ac.jp/dspace/bitstream/2433/275708/1/Sino.Tibetan.lang_5_085.pdf}
\item STEDT classification group page for Southeastern Loloish. This is a classification tree, not a description of Sani. \url{https://stedt.berkeley.edu/~stedt-cgi/rootcanal.pl/group/47/699}
\end{itemize}

\section{Religious and Community Resources}
\begin{itemize}
\item Asia Harvest people group profiles (Sani, and a second profile titled sanie). The Sani profile is the only source in the results that mentions a dictionary: a Sani-Chinese dictionary completed in 1986. I found no title, publisher or other detail. \url{https://people-groups.asiaharvest.org/people-group-profiles/yunnan/sani.pdf} \quad \url{https://people-groups.asiaharvest.org/people-group-profiles/yunnan/sanie.pdf}
\end{itemize}

\section{Digital and Audio Resources}
No audio, video, corpus or software resources specific to Sani were found.

\section{Curriculum Development Context}
Sani has a sizeable speaker base and its own script, but the searches found no teaching materials, textbooks or digital tools. Any curriculum would build on the script and the 1986 dictionary mentioned by Asia Harvest. Those would need to be located and checked first. Script encoding is also unresolved.
\end{document}